\section{Prologue: a shape made of measurements}
\label{sec:prologue}

\subsection{Why look at something finite?}

Physics describes the world with continuous quantities: positions, times, field
strengths, each a real number that can be divided forever. Yet the deepest facts we
know about nature are \emph{discrete}. An electron's spin, measured along any axis,
gives one of exactly two answers. Quarks come in exactly three colours. There are
exactly three generations of matter. The symmetry groups that organise the forces ---
$SU(3)$, $SU(2)$, $U(1)$ --- are continuous, but their representations, which decide
what particles can exist, are labelled by whole numbers.

This programme starts from the discrete end. It asks: is there a single finite object,
small enough to write down completely, whose structure already contains the patterns
physics keeps finding? The candidate studied here is a configuration of forty points
called $\W$. It is not new --- geometers have known it for a century --- but it has an
unusual habit: every time it is examined through a different lens, it turns out to be
a well-known structure from a different part of mathematics or physics. This paper is
an account of those lenses, and of exactly how far each one reaches.

\begin{plainwords}
A useful picture: $\W$ is like a crystal. A crystal is finite in its unit cell, but
the unit cell determines the whole crystal's behaviour --- how it breaks, how light
passes through it, what sounds it makes. The hope of this programme is that $\W$ is a
unit cell for something larger. The honest answer, as you will see, is that it is
a remarkable unit cell for the \emph{symmetries} of physics, and a much weaker one for
its \emph{dynamics}.
\end{plainwords}

\subsection{Bits, trits and qutrits}

Every computer you have used stores information as bits --- switches that are either
off or on. A \emph{trit} is a three-way switch: $0$, $1$, or $2$. On its own that is
not interesting; two bits can imitate a trit. A \emph{qutrit} is different. It is a
physical system with three distinguishable states, obeying quantum mechanics, which
means it can be in a blend of all three settings at once, with relative phases. In
symbols, a qutrit state is a unit vector
\[
\ket{\psi}=a_0\ket{0}+a_1\ket{1}+a_2\ket{2}\in\C^3 ,
\]
and the phases of the complex numbers $a_j$ matter.

Two operations generate everything one can do to the ``shape'' of a qutrit. The
\emph{shift} $X$ moves each setting to the next, and the \emph{clock} $Z$ stamps each
setting with a phase:
\[
X\ket{j}=\ket{j+1 \bmod 3},\qquad Z\ket{j}=\omega^{j}\ket{j},\qquad \omega=e^{2\pi i/3}.
\]
They do not commute: $ZX=\omega\,XZ$. That single twist --- a phase of one third of a
turn picked up when the order of two operations is swapped --- is the seed of
everything in this paper.

\begin{physical}
$X$ and $Z$ are the finite versions of the two most important operations in quantum
mechanics: translation in position and translation in momentum. The relation
$ZX=\omega XZ$ is the three-state version of Heisenberg's $[x,p]=i\hbar$. Physicists
call the group generated by $X$, $Z$ and the phases the \emph{Heisenberg--Weyl group};
its elements are the qutrit \emph{Pauli operators}. They label the possible
measurements one can make on a qutrit that are ``as simple as possible'' --- the
ones every other measurement is built from.
\end{physical}

\subsection{Arithmetic with three numbers}

Because $X^3=Z^3=1$, exponents only matter modulo $3$. We therefore work in the
\emph{field with three elements}, $\F_3=\{0,1,2\}$, where addition and multiplication
are done modulo $3$ --- clock arithmetic on a three-hour clock. It is a genuine
number system: every nonzero element has an inverse ($1\cdot1=1$, $2\cdot2=4=1$).
In particular $2=-1$ and $\tfrac12=2$, facts we will use repeatedly.

A qutrit Pauli operator is named by two numbers, $X^aZ^b$ with $(a,b)\in\F_3^2$.
Two qutrits need four: $(a_1,b_1,a_2,b_2)\in\F_3^4$, a space with $3^4=81$ points.

\subsection{The forty in one breath}

Here is the whole construction, which Sections~\ref{sec:object}--\ref{sec:self} unpack
slowly. Label the $81$ two-qutrit Pauli operators by vectors $x\in\F_3^4$. Remove the
identity, leaving $80$. An operator and its inverse describe the same measurement
(one is the other run backwards), and $x$ and $-x=2x$ label that pair, so the distinct
measurements number
\[
\frac{3^4-1}{3-1}=\frac{80}{2}=40 .
\]
Two of these forty measurements can be performed together exactly when their operators
commute. Recording which pairs commute gives the forty-point geometry $\W$.

\begin{keyidea}
The forty points of $W(3,3)$ are the forty independent directions along which a pair
of qutrits --- or, as Section~3 shows, one qutrit acting on itself --- can be measured.
The lines of the geometry are the complete sets of measurements that can be made at the
same time.
\end{keyidea}

\subsection{What this paper claims, and what it does not}

A ``theory of everything'' should, at a minimum, reproduce the particles and forces of
the Standard Model, the geometry of space-time and gravity, and the numbers --- masses,
couplings, mixing angles --- that we measure. It is important to say at the start that
nothing in this paper achieves that. What it does achieve is a tight web of \emph{exact}
connections between a single finite object and the structures that any such theory
seems to need:
\begin{itemize}
\item the qutrit Heisenberg group and its Clifford symmetries (Sections 2--3);
\item a finite space-time with a light cone, finite clocks, and a rigorous distinction
      between reversible time and an arrow of time (Section 4);
\item the exceptional Lie algebra $E_8$ and its subalgebras $E_7$ and $E_6$ (Section 5);
\item the Freudenthal triple system behind four-dimensional black-hole charges, and a
      $57$-dimensional light cone on which $E_8$ acts by conformal-like maps (Section 6);
\item the octonionic Albert algebra, the Lorentz algebra $\so(1,9)$ and its chiral
      spinor, which carries exactly the quantum numbers of one generation of matter
      (Sections 7--8).
\end{itemize}
Section~9 is devoted to the physics that has been attempted on top of this --- string
compactifications that realise the Standard Model on the same structure --- including
the negative results, and Section~10 to a quantum computer architecture read off the
same geometry. The Epilogue lists what remains open.

\begin{boundary}
The finite geometry supplies \emph{symmetry and structure}. It does not, by itself,
supply a Hamiltonian, an energy scale, a value of Newton's constant, or a reason for any
particular mass. Where this paper uses words like ``time'', ``light cone'' or ``matter'',
it is naming a mathematical structure that behaves like the physical one; whether nature
uses it is a separate, empirical question.
\end{boundary}
